\section{用理想构造域：商环、域扩张与分裂域}


前三节都在讨论"环里的元素怎么分解"。这一节换一个方向：\textbf{用理想构造新的域}。课堂给了两条路线——第一条是把整环"除以"一个极大理想，第二条是取分式域——然后把得到的域用在多项式上，引出分裂域。

\textbf{本节用到的记号：} $R/M$ 为商环；$F,E,K$ 为域；$F[x]$ 为多项式环；$(f)$ 为 $f$ 生成的主理想；$[E:F]$ 为扩域的\textbf{次数}（下文"域扩张与次数"一节定义）；$\operatorname{Frac}(R)$ 为分式域（在下文"分式域：另一条路线"一节构造）；$\mathbb F_p=\mathbb Z/p\mathbb Z$。

\subsection{极大理想给出域，不可约多项式给出极大理想}


\textbf{定理一。} 设 $R$ 是交换含幺环、$M\trianglelefteq R$ 是极大理想，则 $R/M$ 是域。

\textbf{证明。} 由理想对应定理（[Lec 09 §2](../lec09/sections/02-ideals-and-quotient-rings.md)），$R/M$ 的理想与 $R$ 中包含 $M$ 的理想一一对应。$M$ 极大意味着夹在 $M$ 与 $R$ 之间的理想只有 $M$ 和 $R$ 本身，所以 $R/M$ 只有零理想和自身两个理想。而"交换含幺环没有非平凡理想"等价于"它是域"，故 $R/M$ 是域。$\square$

把这条用在多项式环上，就得到一个非常实用的构造。回忆（第 3 节）：$F$ 是域时 $F[x]$ 是欧几里得整环，因而每个理想都是主理想。

\textbf{定理二。} 设 $F$ 是域、$f\in F[x]$ 非零。则

\[
(f)\ \text{是极大理想}\iff f\ \text{是不可约多项式}.
\]

\textbf{证明。} （$\Rightarrow$）反设 $f=gh$，其中 $g,h$ 都不是常数（即都不是单位）。则 $(f)\subsetneq(g)$：包含显然，严格是因为 $\deg g<\deg f$ 意味着 $f$ 不是 $g$ 的常数倍。这与 $(f)$ 极大矛盾。

（$\Leftarrow$）设 $f$ 不可约且 $(f)\subseteq(g)\subseteq F[x]$。因为 $F[x]$ 是 PID，$(g)$ 由某个 $g$ 生成，且 $g\mid f$。由 $f$ 不可约，$g$ 是单位（此时 $(g)=F[x]$）或 $g\approx f$（此时 $(g)=(f)$）。所以中间没有别的理想，$(f)$ 极大。$\square$

\textbf{推论。} $F[x]/(f)$ 是域，当且仅当 $f$ 不可约。

\textbf{例。} 取 $F=\mathbb R$、$f=x^2+1$。它在 $\mathbb R$ 上不可约（判别式 $<0$），所以

\[
\mathbb R[x]/(x^2+1)\cong\mathbb C .
\]

同构把 $x$ 的像送到 $i$。这正是"复数就是给 $\mathbb R$ 添一个 $x^2=-1$ 的解"的严格说法：在商环里 $x^2+1=0$，即 $x$ 扮演了 $i$ 的角色。

\textbf{例。} 取 $F=\mathbb Q$、$f=x^2-2$（在 $\mathbb Q$ 上不可约，因为 $\sqrt2$ 不是有理数），则 $\mathbb Q[x]/(x^2-2)$ 是域，同构于 $\mathbb Q(\sqrt2)=\{a+b\sqrt2:a,b\in\mathbb Q\}$。

\textbf{顺手记下"添一个元素"的记号。} 设 $R'\supseteq R$ 是更大的环、$u\in R'$。记 $R[u]$ 为 \textbf{$R'$ 中包含 $R$ 与 $u$ 的最小子环}，它由 $u$ 的多项式组成：

\[
R[u]=\{r_0+r_1u+r_2u^2+\cdots+r_ku^k:r_i\in R,\ k\ge0\}.
\]

（有限和。）当 $u$ 是某个不可约多项式的根时，$R[u]$ 的加法结构会"降次"——这正是下面 $F[x]/(f)$ 里每个元素都能用 $1,\bar x,\dots,\bar x^{\,n-1}$ 线性表示的原因。分式域、扩域都建立在这个构造之上。

\textbf{例（手写稿给的二次域同构）。} 对无平方因子的整数 $d$，有环同构

\[
\mathbb Z[\sqrt d]\;\cong\;\mathbb Z[x]/(x^2-d).
\]

映射取"把 $x$ 送到 $\sqrt d$"，即 $\varphi\bigl(\sum a_ix^i\bigr)=\sum a_i(\sqrt d)^i$。它满，因为 $\varphi(a+bx)=a+b\sqrt d$ 已取遍 $\mathbb Z[\sqrt d]$；核由带余除法算出：除以首一多项式 $x^2-d$ 得 $f(x)=q(x)(x^2-d)+r(x)$，$\deg r\le1$，于是 $\varphi(f)=0\iff r(\sqrt d)=0$。若 $r=\alpha+\beta x$ 且 $\alpha+\beta\sqrt d=0$，由 $\sqrt d$ 不是有理数得 $\alpha=\beta=0$，故 $\ker\varphi=(x^2-d)$。用第一同构定理即得。

注意这里 $\mathbb Z[x]/(x^2-d)$ 是\textbf{环}但不一定是域（例如 $x^2-10$ 在 $\mathbb Q$ 上不可约，可 $\mathbb Z[x]/(x^2-10)$ 不是域，因为 $\mathbb Z$ 不是域）。要拿到域，系数环必须是域：$\mathbb Q[x]/(x^2-d)$ 才是域。

\subsection{分式域：另一条路线（回顾）}


如果起点是\textbf{整环}而不是域上的多项式环，可以用分式域把除法"补"进去。$R$ 为整环时，在 $R\times(R\setminus\{0\})$ 上定义

\[
(a,b)\sim(c,d)\iff ad=bc,
\]

把等价类记为 $a/b$，按 $\frac ab+\frac cd=\frac{ad+bc}{bd}$、$\frac ab\cdot\frac cd=\frac{ac}{bd}$ 运算，得到\textbf{分式域} $\operatorname{Frac}(R)$。它满足泛性质：任何把 $R$ 单射进域 $K$ 的同态 $\varphi$，都能唯一地延拓为 $\operatorname{Frac}(R)\to K$ 的同态，取值为

\[
\widetilde{\varphi}(a/b)=\varphi(a)\varphi(b)^{-1}.
\]

用一句话说：\textbf{$\operatorname{Frac}(R)$ 是包含 $R$ 的最小域。} 经典例子是 $\operatorname{Frac}(\mathbb Z)=\mathbb Q$ 与 $\operatorname{Frac}(F[x])=F(x)$（有理函数域）。详细证明见 [Lec 09 §4](../lec09/sections/04-units-fractions-characteristic.md)。

两条路线的分工：\textbf{$F[x]/(f)$ 用来"添一个根"，$\operatorname{Frac}(R)$ 用来"添分母"。} 前者要求 $f$ 不可约，后者要求 $R$ 是整环。

\subsection{域扩张与次数}


若域 $E$ 包含域 $F$（$F\subseteq E$ 且两者的加法乘法一致），就说 $E$ 是 $F$ 的\textbf{扩域}，记作 $E/F$。此时有一个关键观察：

\textbf{$E$ 可以看作 $F$ 上的线性空间。} 向量加法就是 $E$ 的加法；标量乘法取 $F$ 中元素去乘 $E$ 中元素——因为 $F\subseteq E$，这个乘法本来就有，且满足线性空间的所有公理。这个线性空间的维数记作

\[
[E:F]=\dim_F E,
\]

称为扩域的\textbf{次数}。$[E:F]<\infty$ 时说 $E/F$ 是\textbf{有限扩张}。

\textbf{例。} $[\mathbb C:\mathbb R]=2$，因为 $\{1,i\}$ 是 $\mathbb R$-基。$[\mathbb Q(\sqrt2):\mathbb Q]=2$，基为 $\{1,\sqrt2\}$。更一般地，若 $f$ 是 $F$ 上 $n$ 次不可约多项式，则 $E=F[x]/(f)$ 满足 $[E:F]=n$，基为 $\{1,\bar x,\dots,\bar x^{n-1}\}$（$\bar x$ 是 $x$ 在商环中的像）。

\textbf{这解释了 4.1 的构造为什么好用：} 想造一个 $F$ 的 $n$ 次扩域，只要找一个 $F$ 上的 $n$ 次不可约多项式，取 $F[x]/(f)$ 即可。

\textbf{顺带一个计数结论。} $E=F[x]/(f)$ 中每个元素都有唯一的代表元 $\bar a_0+\bar a_1\bar x+\cdots+\bar a_{n-1}\bar x^{\,n-1}$（用带余除法把次数降到 $<n$），所以

\[
|E|=|F|^{\deg f}.
\]

当 $F=\mathbb F_p$、$\deg f=n$ 时就得到 $|E|=p^n$——这正是构造 $p^n$ 元有限域的思路：\textbf{找一个 $\mathbb F_p$ 上的 $n$ 次不可约多项式}。

\subsection{特征与素子域}


域的特征定义为

\[
\operatorname{char}(F)=\min\{n\ge1:\underbrace{1+\cdots+1}_{n}=0\},
\]

若不存在这样的 $n$ 则规定为 $0$。因为域是整环，\textbf{特征不是 $0$ 就一定是素数}（[Lec 09 §4](../lec09/sections/04-units-fractions-characteristic.md)）。

考察所有"若干个 $1$ 相加或相减"得到的元素，它们构成 $F$ 的一个子环：特征为 $p$ 时这个子环同构于 $\mathbb Z/p\mathbb Z=\mathbb F_p$；特征为 $0$ 时同构于 $\mathbb Z$。当 $F$ 是域时，特征 $0$ 的情形还可以再取分式域，于是 $F$ 含有一个同构于 $\mathbb Q$ 的子域（因为 $\mathbb Z$ 的分式域是 $\mathbb Q$）。总结：

\[
\text{特征为 }p:\ F\ \text{含子域}\ \mathbb F_p;\qquad
\text{特征为 }0:\ F\ \text{含子域}\ \mathbb Q .
\]

这个\textbf{最小的子域}（任何子域都必须包含它）称为 $F$ 的\textbf{素子域}；[DF] 把它定义为"由 $1_F$ 生成的子域"，中文教材则按"最小性"把它称为\textbf{素域}（$\mathbb Z_p$ 与 $\mathbb Q$ 统称素域）。两种说法指的是同一个对象。

\textbf{有限域的大小。} 设 $F$ 是有限域，则它的特征不能是 $0$（否则含 $\mathbb Q$，已经无限），所以 $\operatorname{char}(F)=p$ 为素数，且 $F$ 含有素子域 $\mathbb F_p$。于是在 4.3 的意义下 $F$ 是 $\mathbb F_p$ 上的有限维线性空间；设 $[F:\mathbb F_p]=n$，则 $F$ 作为 $\mathbb F_p$-向量空间有 $n$ 个基向量，每个基向量的系数有 $p$ 种选法，因此

\[
|F|=p^{\,n}.
\]

\textbf{结论：有限域的阶必为素数幂。} 反过来的问题——对每个素数幂 $p^n$ 是否都存在 $p^n$ 元域，且同构意义下唯一——正是下一节的内容。

\subsection{分裂域}


设 $f\in F[x]$ 是正次数多项式。若扩域 $E\supseteq F$ 满足

\[
f(x)=c\,(x-\alpha_1)(x-\alpha_2)\cdots(x-\alpha_n),\qquad c\in F^\times,\ \alpha_i\in E,
\]

就说 $f$ 在 $E$ 上\textbf{完全分裂}（split），并把每个 $\alpha_i$ 称为 $f$ 在 $E$ 中的根。使 $f$ 完全分裂的"最小"扩域称为 $f$ 的\textbf{分裂域}（splitting field），即 $E=F(\alpha_1,\dots,\alpha_n)$——由 $F$ 和 $f$ 的全部根一起生成的域。

\textbf{"最小"是定义的一部分，不能漏。} 准确的说法是：$f$ 在 $E$ 上完全分裂，且在\textbf{任何}包含 $F$ 的真子域上都\textbf{不}完全分裂。只要求"$f$ 在 $E$ 上分裂"是不够的（把 $E$ 换成更大的域仍然分裂）。

\textbf{次数上界。} 若 $\deg f=n$，则 $[E:F]\le n!$：每添一个根，扩域次数至多等于当前多项式的次数（4.1 的构造给出 $[F_1:F]=\deg g\le n$），而根的总数不超过 $n$，逐次相乘即得。

\textbf{先记住一个反复用到的等价：根 $\iff$ 一次因子。} 设 $E\supseteq F$ 是扩域、$\alpha\in E$，则

\[
f(\alpha)=0\iff(x-\alpha)\mid f(x)\ \text{在}\ E[x]\ \text{中}.
\]

（$\Leftarrow$ 显然；$\Rightarrow$ 用带余除法除以 $x-\alpha$，余式是常数且等于 $f(\alpha)$。）所以"$f$ 在 $E$ 上完全分裂"等价于"$f$ 在 $E[x]$ 中分解成一次因子之积"，也等价于"$f$ 的\textbf{全部}根都在 $E$ 中"。手写稿把这件事写成"$f$ 可解 $\iff$ $f$ 在 $E$ 中可约（能分解）"。

\textbf{定理（存在性）。} 对任意域 $F$ 与任意正次数多项式 $f\in F[x]$，$f$ 的分裂域存在。

\textbf{证明提纲（对 $\deg f$ 归纳）。} 若 $f$ 的所有不可约因子都是一次，则 $f$ 已在 $F$ 上分裂，$E=F$。否则取 $f$ 的一个次数 $\ge2$ 的不可约因子 $g$，则 $F_1=F[x]/(g)$ 是域（定理二）且 $[F_1:F]=\deg g\ge2$；在 $F_1$ 中 $g$ 有一个根 $\bar x$，所以 $f$ 至少多出一个根。把 $f$ 看作 $F_1[x]$ 中的多项式，它在 $F_1$ 上的不可约因子次数严格变小，于是对 $F_1$ 与 $f$ 用归纳假设，得到 $f$ 在某个 $E\supseteq F_1$ 上完全分裂。$\square$

\textbf{"$\bar x$ 是 $g$ 的根"这句话的验证。} 设 $g(x)=\sum a_ix^i$，$\bar x$ 是 $x$ 在 $F_1=F[x]/(g)$ 中的像。在商环里计算 $g(\bar x)=\sum a_i\bar x^{\,i}=\overline{g(x)}=\bar 0$，因为 $g\in(g)$ 在商环中正是零元。所以 $F_1$ 中确实有了 $g$ 的一个根，从而 $f$ 在这个扩域里多分裂出一个一次因子。

\textbf{唯一性证明的关键一步。} 设 $E,E'$ 都含 $F$ 且 $f$ 在两者中都完全分裂。取 $f$ 的根 $u\in E$，它在 $E'$ 中的对应根为 $u'$。则由 $F$ 与 $u$ 生成的子域 $F(u)$ 到 $F(u')$ 有一个保持 $F$ 不动的同态，把 $u$ 送到 $u'$：因为 $u$ 与 $u'$ 在 $F$ 上有\textbf{同一个极小多项式}（都是 $f$ 的不可约因子）——同态只能把根送到根，而 $F[x]/(p)\to E'$、$\bar x\mapsto u'$ 正是这样的同态。把这个步骤对每个根重复，就得到 $E\to E'$ 的 $F$-同构。$\square$

\textbf{定理（唯一性）。} 分裂域在保持 $F$ 不动的同构意义下唯一：若 $E$ 与 $E'$ 都是 $f$ 在 $F$ 上的分裂域，则存在域同构 $\sigma:E\to E'$ 使 $\sigma|_F=\mathrm{id}_F$。

唯一性要读准：\textbf{"唯一"指的是同构唯一，不是集合相等。} 例如 $f=x^2-2$ 在 $\mathbb Q$ 上，$\mathbb Q(\sqrt2)$ 与 $\mathbb Q(-\sqrt2)$ 作为 $\mathbb Q$ 的子域其实是同一个域；但若取两个"形式不同"的构造，它们之间存在保持 $\mathbb Q$ 不动的同构，把 $\sqrt2$ 送到 $\sqrt2$ 或 $-\sqrt2$。课堂上强调"分类肯定不唯一，但同构意义下唯一"就是这个意思。

\textbf{例。} $f=x^2-2$ 在 $\mathbb Q$ 上的分裂域是 $\mathbb Q(\sqrt2)$，因为两个根 $\pm\sqrt2$ 都在其中生成好了：$(-\sqrt2)=-1\cdot\sqrt2$，而 $-1\in\mathbb Q$。相比之下，$f=x^3-2$ 在 $\mathbb Q$ 上的分裂域比 $\mathbb Q(\sqrt[3]{2})$ 大——后者只含一个实根，另外两个是复根，需要再添 $\omega=e^{2\pi i/3}$。

\textbf{下一节预告。} 有限域的存在唯一性（对每个素数幂 $p^n$ 存在唯一的 $p^n$ 元域 $\mathbb F_{p^n}$）将用分裂域证明：$\mathbb F_{p^n}$ 正是 $x^{p^n}-x$ 在 $\mathbb F_p$ 上的分裂域。

\subsection{参考对照}


\begin{itemize}
\item [DF] §13.1（域扩张与次数）、§13.4（分裂域的存在唯一性）、§14.3（有限域）；§9.3（分式域的一般构造）、§8.2（$F[x]$ 的理想与不可约多项式）。
\item [DS] §10.3（特征、素域、定理 10.3.4 $F[x]/(f)$ 是域 $\iff f$ 不可约、定理 10.3.7 有限域的存在唯一性）。该书的极大/素理想只在习题 10.2\#7，且没有域扩张与分裂域。
\item 课程讲义 §8.1（域扩张与次数 8.1、素子域 8.1、有限域大小为素数幂）、§8.2（分裂域的存在唯一 8.4）、§8.3（有限域 8.5）。
\item 原始材料：../../../课堂手记和录音转写/Lec10/。
\end{itemize}
